\section{Measured comparison}
\label{sec:results}

All figures below are on one Apple M4 (10-core, 16\,GB), non-zero-knowledge.
Both provers are run \emph{multithreaded}, so the prover comparison is
core-for-core: \zinc{} with its \texttt{parallel}\,$+$\,\texttt{simd}
build, and \binius{} with its prover parallelised across cores
(rayon\,$+$\,SIMD). \binius{}'s verifier figures, by contrast, are
single-threaded (\Cref{sec:res-scaling}).
\zinc{} instances are indexed by $\nu = \log_2(\text{trace rows})$ with
$\mathrm{compressions}(\nu) = (2^{\nu}-4)/68$, giving $\nu{=}16
\Rightarrow 963$ and $\nu{=}20 \Rightarrow 15{,}420$ compressions;
\binius{} instances are indexed by SHA-256 block count, matched to
\zinc{}'s compression count (block $\approx$ $2^{20}$ at $\nu{=}16$ and
$2^{24}$ at $\nu{=}20$).

\subsection{Headline: the four numbers}
\label{sec:res-headline}

\begin{table}[t]
\centering
\renewcommand{\arraystretch}{1.2}
\begin{tabular}{@{}llrrr@{}}
\toprule
& & \binius{} & \zinc{} & winner \\
\midrule
\multirow{3}{*}{$\approx$\,963 blk ($2^{20}$)}
  & Prove   & $87.5$\,ms  & $98.4$\,ms  & \binius{} $1.13\times$ \\
  & Verify  & $25.9$\,ms  & $4.46$\,ms  & \zinc{} $5.8\times$ \\
  & Proof   & $304$\,KiB  & $1.12$\,MiB & \binius{} $3.8\times$ \\
\midrule
\multirow{4}{*}{$15{,}420$ blk ($2^{24}$)}
  & Prove   & $2.51$\,s   & $1.277$\,s  & \zinc{} $1.97\times^{\ast}$ \\
  & Verify  & $2.39$\,s   & $31.0$\,ms  & \zinc{} $77\times$ \\
  & Proof   & $457$\,KiB  & $2.57$\,MiB & \binius{} $5.6\times$ \\
  & Peak mem& $27.2$\,GB  & $\sim$$13$\,GB & \zinc{} $\sim$$2\times$ \\
\bottomrule
\end{tabular}
\caption{Matched-instance comparison. $^{\ast}$The $\nu{=}20$ prover win
is substantially a \emph{memory} win: \binius{} at $2^{24}$ needs
$\sim$$27$\,GB and is memory-starved on a $16$\,GB host (super-linear);
extrapolating its clean $0.090$\,ms/blk rate to adequate RAM gives
$\sim$$1.4$\,s, i.e. a near-tie. The honest prover verdict is a tie that
\zinc{} wins on memory-constrained hardware.}
\label{tab:headline}
\end{table}

\Cref{tab:headline} is the comparison in four numbers.\footnote{These
\zinc{} prover figures use the Metal GPU commit path and a fresh
same-session sweep, so the $\nu{=}16$ prove ($98.4$\,ms) is the GPU-commit
number; the CPU-commit oblong prove reported in \Cref{sec:impl} is
$\approx 136$\,ms. The $\nu{=}16$ verifier ($4.46$\,ms) is a fresher
measurement that refines the conservative $\approx 11$\,ms quoted in
\Cref{sec:impl}; both are dominated by the same $t$-column proximity test.}
The prover is a
tie ($\binius{}$ $1.13\times$ at $\sim$$950$ blocks; \zinc{} ahead at
$2^{24}$ but mostly on footprint). The two structural gaps are
\zinc{}'s verifier ($5.8\times$, widening to $77\times$) and \binius{}'s
proof ($3.8\times$, widening to $5.6\times$); both widen with $N$ because
one is $\sqrt N$ against $O(N)$ and the other is $\polylog$ against
$\sqrt N$. \zinc{}'s full curve for reference: $\nu{=}9$
$8.03$\,ms\,/\,$1.34$\,ms\,/\,$802$\,KiB; $\nu{=}16$
$98.4$\,ms\,/\,$4.46$\,ms\,/\,$1.12$\,MiB; $\nu{=}20$
$1.277$\,s\,/\,$31.0$\,ms\,/\,$2.57$\,MiB. \binius{}'s curve:
$944$\,blk $87.6$\,/\,$27.1$\,ms\,/\,$304$\,KiB; $4016$\,blk
$362.9$\,/\,$354.7$\,ms\,/\,$404$\,KiB; $15{,}420$\,blk
$2.51$\,s\,/\,$2.39$\,s\,/\,$457$\,KiB.

\subsection{Per-phase decomposition of the prover}
\label{sec:res-phases}

The prover tie is a \emph{net} of large, offsetting per-phase
differences, not a phase-by-phase tie. \Cref{tab:phases} decomposes both
provers at the matched $\sim$$950$-block instance.

\begin{table}[t]
\centering
\renewcommand{\arraystretch}{1.2}
\begin{tabular}{@{}lr@{\quad}lr@{}}
\toprule
\multicolumn{2}{c}{\binius{} ($944$\,blk, $87.6$\,ms)} &
\multicolumn{2}{c}{\zinc{} ($\nu{=}16$, $98.4$\,ms)} \\
\midrule
Commit (RS encode\,+\,Merkle) & $7.1$  & Commit$^{\dagger}$       & $28.0$ \\
IntMul (no-op for SHA)        & $1.0$  & Discharge (oblong $\AND$) & $32.1$ \\
BitAnd reduction              & $24.9$ & UAIR-linear              & $15.4$ \\
Shift reduction               & $43.3$ & Multipoint-eval          & $11.2$ \\
\quad ($\hookrightarrow$ wiring scatter & $27.3$) & Open (PCS)    & $7.6$  \\
PCS open (FRI/BaseFold)       & $10.3$ & z-block                  & $2.8$  \\
\bottomrule
\end{tabular}
\caption{Per-phase prover times (ms). $^{\dagger}$The \zinc{} commit
figure is the \emph{after-prove / repeated-proving} cost; the warm
commit floor is $\sim$$15.5$\,ms --- the $28$\,ms reflects a $\sim$$10$\,ms
cold-cache penalty when the commit fires after a full prove
(\Cref{sec:res-commit}). The single largest \binius{} phase is the Shift
reduction ($49\%$), of which the wiring scatter (\texttt{build\_g\_parts})
alone is $31\%$.}
\label{tab:phases}
\end{table}

Reading \Cref{tab:phases} bucket by bucket:
\begin{itemize}
  \item \textbf{Nonlinearity discharge.} \zinc{}'s oblong discharge
  ($32.1$\,ms) is about $2.1\times$ cheaper than \binius{}'s BitAnd $+$
  Shift reductions ($68.2$\,ms). The dominant \binius{} phase is the
  \emph{Shift reduction} ($43.3$\,ms, $49\%$) --- precisely the
  entry-wise-shift handling that \zinc{} gets for free via $\psi$
  read-offs (\Cref{sec:shifts}). This is the per-phase signature of the
  representational advantage.
  \item \textbf{Commit.} \binius{}'s commit ($7.1$\,ms) is cheaper than
  \zinc{}'s ($28.0$\,ms after-prove; $\sim$$15.5$\,ms warm) --- \zinc{}'s
  rate-$1/4$ codeword (vs \binius{}'s rate-$1/2$) and bit-sliced layout
  cost more even on GPU, and the linear encoder does not fully
  compensate. This is the phase that gives \binius{} its moderate-$N$
  prover edge.
  \item \textbf{Open.} \zinc{}'s flat open ($7.6$\,ms) is \emph{cheaper}
  in time than \binius{}'s FRI open ($10.3$\,ms), even though \zinc{}'s
  opening is far larger in bytes --- the flat-versus-recursive tradeoff
  is time-versus-size, and the proof-size cost of \zinc{}'s opening is not
  a prover-time cost.
  \item \textbf{Plumbing.} \zinc{} carries $\sim$$29$\,ms of dual-projection
  machinery (UAIR-linear, multipoint-eval, z-block) for which \binius{}
  has no analogue (its IntMul is a $1$\,ms no-op).
\end{itemize}
Bucketed, \binius{} wins the commit phase by $\sim$$21$\,ms; \zinc{} wins
the proving core and the open by $\sim$$10$\,ms together; net \binius{}
$+11$\,ms, the measured gap.

\subsection{The verifier and proof scale apart}
\label{sec:res-scaling}

\binius{}'s verifier is essentially the in-the-clear monster sweep
(\Cref{sec:shifts}): measured at $92\%$, $97.5\%$, $99.6\%$ of verify time
at $2^{18}$, $2^{20}$, $2^{22}$, going $6.3 \to 26.5 \to 392$\,ms. The
$26.5\to392$\,ms jump (a $\sim$$15\times$ growth for $4\times$ work) is
super-linear: single-threaded, and at $2^{22}$ the constraint-index
tensor and the multi-million-constraint sweep blow past cache. \zinc{}'s
verifier is $\sqrt N$-bound by its $t$-column proximity test
($4.46 \to 31.0$\,ms from $\nu{=}16$ to $\nu{=}20$, $\sim$$6.9\times$ for
$16\times$ rows). The result is a verifier ratio that grows from
$5.8\times$ to $77\times$. The proof moves the other way: \binius{}'s
FRI proof is $\polylog$ ($304\to404\to457$\,KiB) against \zinc{}'s
$\sqrt N$ flat proof ($0.80\to1.12\to2.57$\,MiB), so the proof ratio
grows from $3.8\times$ to $5.6\times$.

\subsection{A correction the measurement forced}
\label{sec:res-commit}

A careful re-measurement of the \zinc{} commit, isolated, gives a
\emph{warm} floor of $\sim$$15.5$\,ms at $\nu{=}16$ --- the
$\sim$$28$\,ms figure that appears in the per-phase profile is the
after-prove commit, inflated by a $\sim$$10$\,ms cold-cache penalty (the
prior prove's $\sim$$32$\,MB proof and discharge working set evict the
commit's encode/scatter buffers). We verified that this penalty is
\emph{not} GPU-dispatch warmth (a warm-up dispatch yields no compensating
speedup) and \emph{not} reallocation (the commit slab is process-cached);
it is cold memory. The practical consequences: \zinc{}'s warm commit is
$\sim$$2.2\times$ \binius{}'s $7.1$\,ms (not the $4\times$ the raw profile
suggests), and the $\sim$$10$\,ms penalty falls mostly on
\emph{repeated/throughput proving} (and on the benchmark loop, where each
prove follows the last) rather than on a single cold proof.
